% =====================================================================
% Beamer+ paper -- updated to match the current code (2026-10-31).
% Changed passages are marked  % >>> CHANGED  /  % >>> NEW  /  % >>> TODO.
% Unmarked sections are as you had them.
% =====================================================================
\documentclass[sigconf]{acmart}

\usepackage{subcaption}
\usepackage{tikz}
\usepackage{listings}
\usepackage{pdfpages}
\usepackage{minted}

\AtBeginDocument{%
  \providecommand\BibTeX{{%
    Bib\TeX}}}

\setcopyright{acmlicensed}
\copyrightyear{2026}
\acmYear{2026}
\acmDOI{XXXXXXX.XXXXXXX}
\acmConference[Conference acronym 'XX]{Make sure to enter the correct
  conference title from your rights confirmation email}{June 03--05,
  2018}{Woodstock, NY}
\acmISBN{978-1-4503-XXXX-X/2018/06}

\newcommand{\callout}[1]{%
    \tikz[baseline=(n.base)]{%
        \node[circle, fill=darkgray, draw=black, line width=0.4pt,
              text=white, inner sep=1.5pt, font=\footnotesize\bfseries] (n) {#1};%
    }%
}

\begin{document}

\title{Beamer+: An Interactive Presentation Tool}

\author{Chandra Gummaluru}
\email{chandra@cs.toronto.edu}
\affiliation{%
  \institution{University of Toronto}
  \city{Toronto}
  \state{Ontario}
  \country{Canada}
}

\author{Xiaowu Wu}
\email{mika.wu@mail.utoronto.ca}
\affiliation{%
  \institution{University of Toronto}
  \city{Toronto}
  \state{Ontario}
  \country{Canada}
}

\author{Yanting Fan}
\email{yanting.fan@mail.utoronto.ca}
\affiliation{%
  \institution{University of Toronto}
  \city{Toronto}
  \state{Ontario}
  \country{Canada}
}

\author{Hasan Malik}
\email{hasan.malik@mail.utoronto.ca}
\affiliation{%
 \institution{University of Toronto}
 \city{Toronto}
 \state{Ontario}
 \country{Canada}}

\author{Chenhao Zhao}
\email{chenhao.zhao@mail.utoronto.ca}
\affiliation{%
  \institution{University of Toronto}
  \city{Toronto}
  \state{Ontario}
  \country{Canada}}

\renewcommand{\shortauthors}{Gummaluru et al.}

% >>> CHANGED: pens sentence; split widgets; reveal steps.
\begin{abstract}
\LaTeX{} Beamer is widely used for technical presentations, but its static output limits its use in interactive teaching, where an instructor frequently leaves the slides for separate tools such as simulations, coding environments, and polling systems. These switches interrupt the flow of a lecture and cost the attention of everyone in the room. We present Beamer+, a web-based platform that augments existing \LaTeX{} Beamer slide decks with interactive content. A presenter can embed videos, audio, 3D models, and interactive widgets directly onto slides, reveal parts of a slide step by step, annotate in real time with a pen, highlighter, and text tool, and use split-view to display two slides at once, with a single widget able to span both, all without leaving the presentation. Custom widgets are single self-contained HTML files that can be authored without programming, including with the assistance of large language models, so an instructor can build a tool tailored to their own material. We report on two terms of use across three undergraduate computer science courses, describing how Beamer+ was used to teach topics that are difficult to convey with static slides, such as graph search and digital logic, and reflecting on what worked and what did not. We encountered a limitation in authoring structured content during a lecture, where the interface for reshaping a graph on the fly could not keep pace with a live class. We also observed an unanticipated benefit, as students adopted a demonstration widget on their own devices to generate practice problems we had never assigned. We report these experiences in enough detail for other instructors to judge where the tool might fit their own teaching.
\end{abstract}


\begin{CCSXML}
<ccs2012>
 <concept>
  <concept_id>10010405.10010489.10010493</concept_id>
  <concept_desc>Applied computing~Interactive learning environments</concept_desc>
  <concept_significance>500</concept_significance>
 </concept>
 <concept>
  <concept_id>10010405.10010489.10010490</concept_id>
  <concept_desc>Applied computing~Computer-assisted instruction</concept_desc>
  <concept_significance>300</concept_significance>
 </concept>
 <concept>
  <concept_id>10003120.10003121.10003129</concept_id>
  <concept_desc>Human-centered computing~Interactive systems and tools</concept_desc>
  <concept_significance>300</concept_significance>
 </concept>
 <concept>
  <concept_id>10010405.10010489.10010491</concept_id>
  <concept_desc>Applied computing~E-learning</concept_desc>
  <concept_significance>100</concept_significance>
 </concept>
</ccs2012>
\end{CCSXML}
\ccsdesc[500]{Applied computing~Interactive learning environments}
\ccsdesc[300]{Applied computing~Computer-assisted instruction}
\ccsdesc[300]{Human-centered computing~Interactive systems and tools}
\ccsdesc[100]{Applied computing~E-learning}

\keywords{presentation software, interactive learning environments,
  audience response systems, educational technology, lecture tools,
  multimedia authoring, CS education, widgets}

% >>> TODO: retake bp-cover.png -- the toolbar, buttons, and top-right
%           controls have all been redesigned since this screenshot.
\begin{teaserfigure}
  \includegraphics[width=\textwidth]{bp-cover.png}
  \caption{The Beamer+ user interface during a presentation.}
  \Description{The Beamer+ interface showing a slide with an embedded
  interactive widget, the slide navigator, and the annotation toolbar.}
  \label{fig:teaser}
\end{teaserfigure}

\maketitle

\section{Introduction}
\LaTeX{} Beamer has long been the standard tool for creating technical presentations, owing to its philosophy of separating content from design, native support for mathematical and scientific notation, and the ease with which both the source files and output can be shared. However, a significant drawback of the system is that it can only create static presentations, offering limited support for multimedia or interactive elements. While other presentation tools offer some support for multimedia, they lack the qualities that make \LaTeX\ Beamer well suited to technical presentations. They also provide little support for custom interactive elements.

% >>> CHANGED: feature list (pens, reveal steps, split widgets, recording/notes).
To address these limitations, we developed Beamer+, a web-based platform that augments static slide decks with multimedia and interactive elements. A presenter can embed videos, audio, 3D models, and widgets directly into slides, uncover parts of a slide one step at a time, annotate with a pen, highlighter, and text tool, and use split-view mode to display two slides simultaneously, with a single widget able to span both. A lecture can also be recorded, and saved as a handout of notes that captures the annotations and the final state of each widget. The result is a tool that brings the full range of today's interactive content into a presentation, without giving up the benefits of \LaTeX\ Beamer.

Beamer+ grew out of our own experiences within the classroom. In Section~\ref{sec:motivation}, we describe the specific challenges that prompted its development. In Section~\ref{sec:features}, we present the features of Beamer+ that address these challenges. In Section~\ref{sec:experience}, we report on our experience using Beamer+ across several undergraduate computer science courses. Section~\ref{sec:related} discusses related work, and Section~\ref{sec:conclusion} concludes with a discussion of future directions.

\section{Motivation}
\label{sec:motivation}
Beamer+ grew out of a recurring problem encountered while teaching an undergraduate algorithms and data structures course. In preparing and presenting a lecture on Dijkstra's algorithm, it became evident that using static slides to cover such a multi-step process came with some drawbacks. First, the slides were laborious to produce. Conveying a single run of the algorithm meant building a separate figure for each step, resulting in an unwieldy set of nearly identical slides. Assembling these figures by hand was time-consuming, and keeping them consistent was error-prone, since a single change to the example had to be propagated carefully across every slide that followed. Second, the sheer number of slides made the example difficult to follow and hard to annotate. Any marking added to emphasize a particular detail would disappear when advancing to the next slide, since each step was a separate static image rather than a persistent canvas. Finally, and most importantly, the presentation was rigid. The aforementioned slide set showed one fixed run of the algorithm on one fixed graph, so when a student asked how the search would change if an edge weight were modified, there was no practical way to explore the question in the moment. Answering it meant either working through the change by hand on the board or deferring it, neither of which preserved the continuity of the lecture.

\begin{figure}[h]
  \centering
  \begin{subfigure}[t]{0.48\linewidth}
    \includegraphics[page=1, width=\linewidth]{images/dijkstra_slides2.pdf}
  \end{subfigure}
  \hfill
  \begin{subfigure}[t]{0.48\linewidth}
    \includegraphics[page=2, width=\linewidth]{images/dijkstra_slides2.pdf}
  \end{subfigure}

  \vspace{0.5em}

  \begin{subfigure}[t]{0.48\linewidth}
    \includegraphics[page=3, width=\linewidth]{images/dijkstra_slides2.pdf}
  \end{subfigure}
  \hfill
  \begin{subfigure}[t]{0.48\linewidth}
    \includegraphics[page=4, width=\linewidth]{images/dijkstra_slides2.pdf}
  \end{subfigure}
\caption{Four consecutive steps of Dijkstra's algorithm, each differing only in which node is being visited and the tentative distances recorded at each node.}
  \label{fig:dijkstra}
\end{figure}

Each subsequent offering of the course prompted an attempt to do better, but each fell short in an instructive way. Replacing the slides with a pre-rendered animation reduced the size of the deck, but it introduced problems of its own. The animation had to be played in a separate window, and once it was playing there was no natural way to pause and annotate it to discuss a particular step. More importantly, using a pre-rendered animation still meant that there was no easy way to explore students' impromptu questions. Moving the presentation into software with stronger multimedia support made animations easier to embed, but sacrificed the mathematical typesetting and source-based workflow that make \LaTeX\ Beamer well suited to technical material. Rendering Beamer slides as images and importing them into that software restored the typesetting, but made the deck onerous to maintain, since any change meant repeating the entire process.

Across these attempts the underlying problem stayed the same. The interactive content always lived outside the presentation, and uniting the two required either significant manual effort or a separate application. The capability that was missing throughout, the ability to respond to a student's question by exploring it live, is exactly the kind of interactive activity that makes a lecture more than a static recitation. Such activities are not specific to search algorithms. Many technical concepts are best understood by watching a process unfold, through a simulation, a running program, or an interactive visualization.

Interactive activities of this kind are a well-supported part of effective teaching. In a meta-analysis of 225 studies across science, engineering, and mathematics, Freeman et al.\ found that active learning improved examination performance and reduced failure rates relative to traditional lecturing \cite{freeman2014active}. Running a simulation, executing code, polling the room, or working through a problem in a shared workspace are all ways of introducing active learning, but each typically requires its own application, so an instructor must repeatedly switch between the presentation and the tools that support these activities.

Research in cognitive psychology has documented a measurable cost to task switching \cite{monsell2003task}, as attention remains partly fixed on the previous task and lowers performance on the current one \cite{leroy2009attention}. These effects were established with controlled laboratory tasks rather than classroom activities, but an instructor who alternates between explaining material and operating a separate application is engaging in a similar kind of switching, and the resulting loss of attention can diminish the quality of instruction. Students must reorient themselves as well, each time the lecture moves between interfaces. In their review of audience response systems, Kay and LeSage identify the time required to set up and operate such tools as a recurring drawback that can undermine their pedagogical benefit \cite{kay2009audience}, and similar friction can be expected whenever an interactive activity depends on software outside the presentation.

Beamer+ was created to remove this separation. By bringing interactive content directly into \LaTeX\ Beamer presentations, it lets an instructor incorporate simulations, executable code, audience interaction, and other dynamic elements without leaving the slides. The aim is not merely to add multimedia to Beamer, but to make interactive teaching activities a natural part of the presentation itself.

\section{Features}
\label{sec:features}
% >>> CHANGED: opening/saving (no "upload/export"), sessions, what Save writes.
A Beamer+ presentation is a \texttt{.zip} file containing a static \texttt{.pdf} slide deck along with any overlays, annotations, and metadata. A presenter starts a session, opens either a plain \texttt{.pdf} or a saved \texttt{.zip}, and augments it with overlays in edit mode (see Section~\ref{sec:edit-mode}). Overlays can be freely positioned and resized on any slide, and since the metadata is stored separately from the \texttt{.pdf} file, the original slide deck remains unchanged. A presenter can also insert blank slides as space for annotations (see Section~\ref{sec:annotations}), create a ``split'' that displays two slides side by side at a designated point in the presentation (see Section~\ref{sec:split-view}), and mark slides as hidden so they are skipped while presenting. When the session is closed, the presentation is saved as a \texttt{.zip} that reopens at the slide where the presenter left off (see Section~\ref{sec:sessions}).

% >>> TODO: retake overview.pdf / overview-split.pdf and renumber the
%           callouts. The annotation toolbar is now a rail docked on the
%           right (a bar along the top in split view), and the top-right
%           corner holds Edit and Close Session.
\begin{figure}[H]
    \centering
    \begin{subfigure}[b]{\linewidth}
        \centering
        \includegraphics[width=\linewidth]{images/overview.pdf}
        \caption{Overview}
        \label{fig:overview}
    \end{subfigure}
    \hfill
    \begin{subfigure}[b]{\linewidth}
        \centering
        \includegraphics[width=\linewidth]{images/overview-split.pdf}
        \caption{Overview (split)}
        \label{fig:overview-split}
    \end{subfigure}
    \caption{The Beamer+ interface in (a) single-slide and (b) split-view modes. Numbered callouts are referenced throughout this section.}
    \label{fig:overview-figure}
\end{figure}

% >>> CHANGED: bookmarks added.
\subsection{Navigation}
A presenter can navigate between slides using the left and right arrow keys, or jump directly to any slide by clicking its thumbnail in the \textit{slide navigator} \callout{1}. Hidden slides are skipped when using the arrow keys, making it easy to include optional or supplementary material without disrupting the flow of a presentation. Slides can also be bookmarked; bookmarks appear as pins in the navigator, so a presenter can return to a key slide in one click.

% >>> NEW subsection (reveal steps).
\subsection{Reveal Steps}
\label{sec:reveals}
Beamer's \texttt{\textbackslash pause} uncovers a slide one part at a time, but only by producing one PDF page per step. Beamer+ offers the same effect on a single page: in edit mode, a presenter draws \emph{reveal boxes} over regions of a slide and numbers them. While presenting, the slide begins with every box closed, and each press of the right arrow key uncovers the next step before moving to the next slide; the left arrow key closes it again. A box can either blend into the slide's background, so the hidden region looks empty, or appear as a visible frosted panel. Videos and widgets can likewise be set to appear on a given step. Annotations made beneath a closed box stay hidden until it opens.

% >>> CHANGED: split toggle location; split widgets added.
\subsection{Split-View Mode}
\label{sec:split-view}
A presenter can display two slides simultaneously using \textit{split-view} \callout{2}, toggled from the annotation toolbar. The relative size of each pane can be adjusted by dragging the \textit{divider} \callout{8} between them, and each slide is fitted to its pane so that both remain as large as possible.
% >>> TODO: describe how a slide is assigned to each pane while presenting
%           (your old text said the navigator thumbnails are split into
%           left/right halves -- check this still matches, then add it back).

A split view can also hold a \emph{split widget}: a single widget that spans both slides, with one half on each. The first such widget is the Python workspace, which places its code editor on the left slide and, on the right, a shell, a memory diagram, or a debugger, between which the presenter can switch with tabs during the lecture. Because it is one widget, running code on the left prints its output on the right, and the program's state survives leaving the split and returning to it. When either slide is shown on its own, its half is drawn as it was last left.

% >>> CHANGED: one pen + one highlighter; new tools; toolbar contents.
\subsection{Annotations}
\label{sec:annotations}
The \textit{annotation toolbar} \callout{4} provides a pen and a highlighter, each with five sizes and eight colours, together with a text tool, an eraser with adjustable size, undo and redo, and a button that clears the slide. Drawing a stroke and holding still snaps it to a clean shape. To direct the audience's attention, a spotlight dims the slide around the cursor, and a magnifier enlarges the region beneath it. The toolbar also holds slide actions: toggling split view, inserting a blank page, duplicating the current slide, and hiding it. Annotations are saved with the presentation, so they are preserved across sessions.
% >>> TODO: retake annotation-toolbar-horizontal.pdf (new tool set; it is
%           now a vertical rail in single view).
\begin{figure}[H]
    \centering
    \includegraphics[width=\linewidth]{images/annotation-toolbar-horizontal.pdf}
    \caption{The annotation toolbar, providing a pen, highlighter, text tool, eraser, spotlight, magnifier, undo and redo, and slide actions.}
    \label{fig:annotation-bar}
\end{figure}

\subsection{Overlays}
\label{sec:overlays}
% >>> CHANGED: audio added.
Beamer+ uses the term \emph{overlay} to refer to any interactive or multimedia element layered on top of a slide's static content. A presenter can augment slides with several types of overlay, including videos, audio, 3D models, and interactive widgets \callout{12}. The annotation canvas always sits above overlays, so a presenter can write directly on top of any interactive content, just as they would on a static slide.

% >>> CHANGED: removed "playback position is saved" (not in the current code);
%     play modes; expand-on-play.
\subsubsection{Videos and Audio}
Video overlays embed clips (e.g., demonstrations, recorded experiments, or process animations) directly onto a slide. A clip can play when clicked, start automatically when its slide is reached, or loop, and a video can be set to expand to fill the slide a few seconds after it starts playing. Audio overlays work the same way, for sound clips that need no picture.

% >>> CHANGED: removed "camera viewpoint is saved" (not in the current code).
\subsubsection{3D Models}
Model overlays render interactive 3D objects (e.g., molecular geometries, mechanical assemblies, or anatomical structures) on a slide from \texttt{.glb} files. The model can be rotated, panned, and zoomed during a presentation to examine detail from any angle.

% >>> CHANGED: current built-in list; full vs overlay; slide-native look;
%     editing from the widget itself.
\subsubsection{Widgets}
Widgets are the primary mechanism for making static slides dynamic. Beamer+ includes a built-in library of thirteen widgets: a calculator, a function plotter, a circuit simulator, a Jupyter-style notebook, a Python workspace, a terminal, a timer, a randomizer, a map, a PDF textbook viewer, a YouTube player, a camera feed, and an audience-response widget (Section~\ref{sec:audience}). A widget either fills the whole slide or, where it supports it, sits as a small box over part of a content slide, as a timer or calculator might. The built-in widgets are drawn to match a \LaTeX\ slide, using Computer Modern type and TikZ-like outlines, so they read as part of the slide rather than as an embedded application. Each widget is configured from the widget itself: pressing its \emph{Edit} button switches it into an edit mode in which most widgets are set up in place, so that, for example, a timer's duration is typed into its digits and a poll's question and answers are typed where they will appear. A widget's state is saved with the presentation, so that it resumes where it left off when the presenter returns to the slide.

% >>> CHANGED: schema id/format, config delivery, no Properties panel,
%     no WIDGET_GUIDE.md; inlineEdit / split / print hooks; self-contained.
\subsubsection{Custom Widgets}
In addition to the built-in library, Beamer+ supports custom widgets that allow an instructor or developer to create domain-specific learning tools tailored to their course. A widget is a single self-contained \texttt{.html} file. Its editable settings are declared in a \texttt{JSON} schema block in the file's \texttt{<head>}, giving each field's key, label, type (e.g., text, number, checkbox, or a choice from a list), and default. Beamer+ reads this schema and builds a settings form for the widget automatically, with no additional configuration required; the widget receives its current settings in \texttt{window.WIDGET\_CONFIG} and is sent a \texttt{widget-config} message whenever they change (Listing~\ref{lst:widget-skeleton}). Optional schema entries let a widget opt into further behaviour: \texttt{modes} declares whether it can appear as an overlay, \texttt{inlineEdit} replaces the generated form with the widget's own in-place editing, and \texttt{split} lets it span both slides of a split view. A widget can also register a print handler to control how it appears in the notes handout (Section~\ref{sec:sessions}). Because a widget carries everything it needs in its one file, it also runs when opened directly in a browser, outside Beamer+.
\begin{listing}[H]
\begin{minted}[frame=single, fontsize=\small, linenos]{html}
<!DOCTYPE html>
<html>
<head>
  <meta charset="UTF-8">
  <script id="widget-schema" type="application/json">
    {
      "label": "My Widget",
      "modes": ["full", "overlay"],
      "fields": [
        { "key": "heading", "label": "Heading",
          "type": "text",   "default": "Hello" },
        { "key": "count",   "label": "Count",
          "type": "number", "default": 5 }
      ]
    }
  </script>
</head>
<body>
  <div id="widget-root"></div>
  <script>
    function render(cfg) {
      const root = document.getElementById("widget-root");
      root.innerHTML = `<h2>${cfg.heading ?? "Hello"}</h2>`;
      // Use cfg.count below...
    }
    // Settings when the widget loads...
    render(window.WIDGET_CONFIG || {});
    // ...and every time they change.
    window.addEventListener("message", (e) => {
      if (e.data?.type === "widget-config")
        render(e.data.config);
    });
  </script>
</body>
</html>
\end{minted}
\caption{Minimal custom widget skeleton. The \texttt{widget-schema} block declares editable settings; \texttt{render} draws the widget from their current values, on load and after every change.}
\label{lst:widget-skeleton}
\end{listing}
% >>> TODO: WIDGET_GUIDE.md no longer exists in the repo. Either restore a
%           guide file and keep its name below, or point to the README.
For a presenter without web development experience, modern large language models (LLMs) are well suited to generating custom widgets from a natural-language description. To facilitate this, Beamer+ documents the full widget specification, schema syntax, and worked examples in a single guide that can be provided directly as context when prompting an LLM.

% >>> CHANGED: panel contents; widget settings live in the widget; blank-slide
%     paper; duplicate/delete/bookmark; split ratio; reveal boxes; saving.
\subsection{Edit Mode}
\label{sec:edit-mode}
A presentation can be edited by switching into edit mode. In this mode, selecting a slide brings it into the main editing area, where overlays and reveal boxes can be added and arranged on top of the rendered page. Selecting an overlay opens a properties panel holding its name, text size, position and size, and the step on which it appears; a widget's own settings are edited from the widget, as described above. Slides can be reordered by dragging their thumbnails in the slide navigator, duplicated, deleted, bookmarked, or marked as hidden, which excludes them from the presentation while keeping them in the deck for reference or later use. Blank slides can be inserted at any point as additional working space, with a choice of plain, lined, grid, or dotted paper, and support the same overlays and annotations as other pages. A presenter can also create a ``split,'' which pairs two existing slides at a ratio chosen ahead of time and displays them side by side when the split is reached while presenting; a split can additionally hold a split widget (Section~\ref{sec:split-view}).
% >>> TODO: retake overview-edit.pdf (new panel and toolbar).
\begin{figure}[H]
    \centering
    \includegraphics[width=\linewidth]{images/overview-edit.pdf}
    \caption{The Beamer+ editor. Overlays are added and arranged on the slide, and a properties panel exposes the name, size, and position of the selected overlay.}
    \label{fig:edit-overview}
\end{figure}

% >>> CHANGED: rewritten -- one widget with nine question types.
\subsection{Audience Response}
\label{sec:audience}
Beamer+ includes a single audience-response widget that supports nine question types: multiple choice, true/false, rating scale, numeric estimate, numerical answer, word cloud, fill in the blanks, a two-dimensional plot on which each student places a point, and open-ended text. Questions and answers may contain Markdown, code, and \LaTeX{} mathematics, edited in place on the slide. While a poll is open, the widget shows a QR code that audience members scan to answer from their own devices; each submission is relayed to the presenter via Socket.IO, so the response count updates in real time. A poll can be closed by hand or after a countdown. Results are drawn on the slide in the same plain style as the slides themselves: bar charts for choices and scales, a cloud for words, the placed points with their average for plots, and the responses themselves for blanks and open text. A correct answer, when one is set, can be revealed after the results are shown.

Open-ended responses can additionally be passed to a pluggable AI summarization backend. A summarizer is a Python file, placed either on the server or inside a presentation's \texttt{.zip}, that exports a single function taking the responses and returning a ranked list of summary and respondent-count pairs. Two summarizers, based on commercial language models, are included. The summaries are displayed one at a time on the slide, and the presenter steps through them.

% >>> NEW subsection (replaces "Usage Guide"): sessions, saving, recording, notes.
\subsection{Sessions, Recording, and Notes}
\label{sec:sessions}
Beamer+ runs from a single server, either on the presenter's own machine or hosted by an institution for many concurrent lectures. Each presenter works in a separate session at its own address, and the QR codes shown to students point into that session. A first-time user can choose \emph{Try Demo}, which opens a demonstration presentation with a guided tour of the main features.

A lecture can be recorded from within Beamer+, capturing the slides, annotations, and widgets together with the presenter's microphone and the sound of any videos or widgets. Recordings are written to the browser's storage as they are made, so a reload or crash does not lose them.

When the session is closed, the presenter can save up to three things: the presentation as a \texttt{.zip}, any recordings, and a \emph{notes} PDF intended as a handout. The notes contain every visible slide with its annotations drawn on, each widget as it stood at the end of the lecture, a still frame for each video, and the results of each poll. A widget can add pages of its own after its slide: the circuit simulator prints its board, the notebook prints every cell, and the Python workspace prints its program, shell transcript, and memory diagram. Slides from the original \texttt{.pdf} are copied as vector pages, so their text remains selectable.

\section{Classroom Experience}
\label{sec:experience}
We have used Beamer+ over two terms in three undergraduate computer science courses: an introduction to artificial intelligence, a course on data structures and algorithms, and a course on computer organization. All three were taught by one of the authors, who used Beamer+ for a handful of lectures in each course, focusing on the topics that are hardest to convey with static slides. We did not conduct a controlled study; instead, this section reports on how the tool was used and what we learned from using it.

The search widget that motivated Beamer+ (Section~\ref{sec:motivation}) was used in both the artificial intelligence and the data structures courses, where graph search recurs. Embedded directly on the slide, it animated the progress of an algorithm such as A* and allowed its parameters to be changed during the lecture. When a student asked how the search would behave under a different edge weight, the new value could be entered and the algorithm re-run in front of the class, rather than deferred to office hours or sketched by hand. This was precisely the capability whose absence had originally made the topic awkward to teach.

Using the widget in front of students also revealed a limitation we had underestimated while building it. Editing edge weights on the fly worked well, but changing the structure of the graph itself, adding or removing nodes and edges, was awkward, because the widget had no convenient interface for doing so mid-lecture. A student who asked to see the algorithm run on a differently shaped graph could not always be accommodated in the moment. The tool had been built to escape the rigidity of static slides, yet in this respect it imposed a narrower rigidity of its own. The gap was specifically in the presenter-facing interface for authoring a graph during a lecture, not in the widget's ability to handle different graphs, and addressing it is a clear direction for future work.

That underlying flexibility surfaced in an unexpected way. Although the widget was built as a tool for demonstration at the front of the room, students began using it themselves during and after lecture, generating new graphs on their own devices and running the algorithm on them to practice. A tool intended to support explanation became, with no change on our part, a way for students to create their own practice problems. This was among the more valuable outcomes of the deployment, and one we had neither designed for nor anticipated.

% >>> CHANGED: "word-cloud widget" -> the audience-response widget's word cloud.
In the computer organization course, Beamer+ was used to teach digital logic, where a circuit and its timing behaviour are difficult to present together on static slides. A logic-simulator widget allowed inputs to be toggled live, split-view paired a circuit diagram with its timing diagram, and the annotation canvas was used to draw the correspondence between a change in the circuit and the resulting edge in the waveform. Across all three courses, a word-cloud poll served as a warm-up that surfaced which concepts students found most confusing, and short video overlays and 3D models were embedded where a static figure was insufficient.

Taken together, our experience suggests that Beamer+ achieved its central aim of keeping a lecture's activities within a single environment: demonstration, comparison, annotation, and audience response were carried out without leaving the presentation. The main shortcoming we encountered was the difficulty of authoring structured content on the fly, and the main surprise was the tool's adoption by students for their own practice. Both point toward the same next step: a faster interface for creating and editing interactive content, usable by instructors and students alike.

\section{Related Work}
\label{sec:related}
Slide-based software such as PowerPoint, Keynote, Google Slides, and \LaTeX\ Beamer arranges content as a sequence of slides, and web-based tools such as reveal.js bring the same model into the browser. A body of research has extended this model in several directions. Zoomable user interfaces, popularized by Prezi and investigated in research systems such as CounterPoint \cite{good2002zoomable} and Fly \cite{lichtschlag2009fly}, replace the fixed sequence with a spatial canvas, and pen-based systems such as Classroom Presenter bring handwriting into the lecture, allowing an instructor to ink directly on slides and share digital ink with students \cite{anderson2007classroom}. Closest to our work is MindXpres, a web-based, extensible presentation platform whose plug-in architecture allows interaction components to be added to a presentation rather than chosen from a fixed set \cite{roels2014mindxpres}. Beamer+ differs in two respects that bear on classroom adoption: it augments an existing \texttt{.pdf} deck in place rather than requiring content to be re-authored in a dedicated format, and its custom widgets can be created without programming, including from a natural-language description.

The activities an instructor reaches for outside the deck have each been studied as a medium in their own right. Interactive simulations let students manipulate a model and observe its behaviour directly, as in the PhET project, which reports learning gains when its simulations replace static presentation of the same material \cite{wieman2008phet}. Explorable explanations embed reactive content that lets a reader change an author's assumptions and see the consequences \cite{victor2011explorable}, and authoring tools such as Idyll reduce the effort of building such interactive articles \cite{conlen2018idyll}. Computational notebooks place executable code beside prose and output, descending from Knuth's notion of literate programming \cite{knuth1984literate}, and have been used in the classroom to run and edit code live rather than display a static listing \cite{johnson2020jupyter}. Audience response systems such as Kahoot, Mentimeter, and Poll Everywhere bring student interaction into the lecture, but offer a fixed menu of interaction types. Each of these media is typically delivered by a separate application; Beamer+ brings them into a single presentation surface and additionally allows arbitrary custom interactive content to be embedded alongside them.

\section{Conclusion}
\label{sec:conclusion}
We presented Beamer+, a web-based tool that augments \LaTeX{} Beamer slide decks with interactive overlays, and reported on two terms of use across three undergraduate computer science courses. The design rests on a single premise, that the cost of leaving the slides for a separate application is paid not only in lost time but in the attention of everyone in the room. By bringing simulations, executable code, audience response, and other dynamic content onto the slide itself, Beamer+ allowed topics that are hard to convey with static slides to be taught within one environment, combining demonstration, comparison, annotation, and audience response without leaving the presentation. Custom widgets extend this beyond a fixed menu of interactions, since an instructor can author a tool for their own material as a single HTML file with no web development experience required.

Putting the tool in front of students taught us more than building it. Our experience surfaced a limitation, the difficulty of authoring structured content such as a graph during a lecture, where the demonstration interface could not keep pace with a class that wanted to reshape an example on the spot. It also surfaced an unanticipated benefit, the adoption of an instructor's demonstration widget by students, who used it on their own devices to generate practice problems we had never assigned. A tool built for the front of the room found a second use in the hands of the people it was meant to reach.

% >>> TODO: replace the placeholder URL.
These observations shape our future work. The most pressing improvement is a faster interface for creating and editing structured interactive content on the fly, which would address the limitation we encountered and support the student-driven practice we observed. We also plan to broaden deployment beyond a single instructor and to study more systematically how the tool affects teaching and learning. By describing both what worked and what did not in concrete terms, we hope to give other instructors a clear sense of where Beamer+ might serve their own classrooms. Beamer+ and its widget specification are available at \url{https://example.com/beamerplus}.

\begin{acks}
Removed for review.
\end{acks}

\bibliographystyle{ACM-Reference-Format}
\bibliography{sample-base}

\end{document}
